\section{Reconstruction tomographique parcimonieuse par ADMM et ondelettes Bior4.4}
%------------------------------------------------
\begin{frame}{Résultats Quantitatifs}
    \only<1>{
        \begin{block}{Performance de la méthode ADMM-Bior4.4}
            \textbf{Cas d'étude détaillé (Image Shepp-Logan)}
            \begin{table}[H]
            \centering
            \caption{Évaluation comparative des performances de reconstruction}
            \label{tab:quantitative_results}
            \begin{tabular}{lcc}
            \hline
            \textbf{Méthode} & \textbf{PSNR $\uparrow$ (dB)} & \textbf{SSIM $\uparrow$}  \\
            \hline
            FBP & 25.04 & 0.678 \\
            ISTA & 24.83 & 0.768  \\
            FISTA & 31.50 & 0.758  \\
            ADMM-Bior4.4 & \textbf{37.31} & \textbf{0.972}  \\
            \hline
            \end{tabular}
            \end{table}
            
            \begin{itemize}
                \item Ces valeurs, obtenues sur l'image de référence, démontrent une reconstruction de très haute qualité, avec une excellente préservation des structures anatomiques.
            \end{itemize}
        \end{block}
    }
    % \only<2>{
    %     \begin{block}{Performances moyennes sur l'ensemble des 19 images}
    %         \begin{table}[h]
    %         \centering
    %         \caption{Comparaison des métriques PSNR et SSIM}
    %         \begin{tabular}{|c|c|c|c|c|}
    %         \hline
    %         \textbf{Algorithme} & \textbf{PSNR\_mean} & \textbf{PSNR\_std} & \textbf{SSIM\_mean} & \textbf{SSIM\_std} \\
    %         \hline
    %         ADMM-Bior4.4 & 26.203727 & 0.151281 & 0.865829 & 0.008924 \\
    %         FISTA         & 25.132216 & 0.220781 & 0.677443 & 0.010547 \\
    %         ISTA          & 22.397596 & 0.385292 & 0.650520 & 0.015583 \\
    %         FBP           & 21.610267 & 0.301994 & 0.558846 & 0.021965 \\
    %         \hline
    %         \end{tabular}
    %         \end{table}
    %     \end{block}
    % }
    \only<2>{
        \begin{block}{Performances moyennes sur l'ensemble des 19 images}
            \centering
            \captionof{table}{Comparaison des métriques PSNR et SSIM}
            \resizebox{0.9\textwidth}{!}{
                \begin{tabular}{|c|c|c|c|c|}
                \hline
                \textbf{Algorithme} & \textbf{PSNR\_mean} & \textbf{PSNR\_std} & \textbf{SSIM\_mean} & \textbf{SSIM\_std} \\
                \hline
                ADMM-Bior4.4 & 26.203727 & 0.151281 & 0.865829 & 0.008924 \\
                FISTA         & 25.132216 & 0.220781 & 0.677443 & 0.010547 \\
                ISTA          & 22.397596 & 0.385292 & 0.650520 & 0.015583 \\
                FBP           & 21.610267 & 0.301994 & 0.558846 & 0.021965 \\
                \hline
                \end{tabular}%
            }
        \end{block}
    }
    \only<3>{
        \begin{block}{Analyse et Interprétation}
            \begin{itemize}
                \item \textbf{Analyse :} ADMM-Bior4.4 surpasse nettement les autres méthodes en moyenne, avec des écarts-types plus faibles, attestant d'une grande stabilité.
                
                \vspace{0.3cm}
                
                \item \textbf{Interprétation :}
                \begin{itemize}
                    \item Un \textbf{PSNR} de 26.20 dB indique une excellente fidélité aux données.
                    \item Un \textbf{SSIM} de 0.866, proche de 1, confirme une très bonne préservation des structures anatomiques, même en contexte de sous-échantillonnage.
                \end{itemize}
            \end{itemize}
        \end{block}
    }
\end{frame}